\documentclass{article}
\usepackage{amsmath, amssymb}
\usepackage{enumitem}

\title{Learning \LaTeX\ Through a Simple Article}
\author{Kotapati Sai Mounika B24CS019}
\date{04/09/2025}

\begin{document}
	\maketitle
	
	\section{Introduction}
	This article is designed to show how a complete \LaTeX\ document looks. It
	demonstrates document classes, titles, sectioning, lists, and mathematical
	symbols.
	
	\subsection{Motivation}
	\LaTeX\ is used for producing professional documents.
	
	\subsubsection{Why Students Need It}
	Because it ensures formatting is consistent.
	
	\textbf{Key Point} Learning \LaTeX\ early saves time in research.
	
	\textbf{Mini Note} Even small reports benefit from \LaTeX.
	
	\section{Organizing Information}
	Information is often presented as lists. \LaTeX\ supports both itemized and
	numbered lists.
	
	\subsection{Custom Itemize}
	\begin{itemize}[label=$\bullet$]
		\item Normal bullet
		\item[\!] Urgent task
		\item[$\blacksquare$] Completed work
	\end{itemize}
	
	\textbf{NOTE} Special note
	
	\subsection{Numbered List}
	\begin{enumerate}
		\item First numbered point
		\item[\!] Custom exclamation
		\item[$\blacksquare$] Black square instead of number
		\item[NOTE] Replace number with text
	\end{enumerate}
	
	\subsection{Nested Example}
	\begin{enumerate}
		\item Outer Level
		\begin{itemize}
			\item Inner bullet
		\end{itemize}
	\end{enumerate}
	
	\textbf{Note:} This is a note
	
	\textbf{Wait!} This is a warning
	
	\section{Mathematics in Action}
	By practicing this, you now cover all the essential commands from the tutorial.
	
	\subsection{Summation Example}
	\[
	S_n = \frac{1}{n} \sum_{k=1}^{n} \frac{k^2 + k + 1}{k^2 + 1}
	\]
	
	\subsection{Geometric Series}
	\[
	1 + r + r^2 + r^3 + \cdots + r^n = \frac{1 - r^{n+1}}{1-r}
	\]
	
\end{document}